\documentclass[11pt,a4paper]{article}

% ---- Packages ----
\usepackage[utf8]{inputenc}
\usepackage[T1]{fontenc}
\usepackage{amsmath,amssymb,amsthm}
\usepackage{mathtools}
\usepackage[margin=1in,headheight=14pt]{geometry}
\usepackage{hyperref}
\usepackage{enumitem}
\usepackage{xcolor}
\usepackage{tcolorbox}
\usepackage{fancyhdr}
\usepackage{booktabs}

% ---- Hyperref ----
\hypersetup{
    colorlinks=true,
    linkcolor=red,
    citecolor=blue,
    urlcolor=blue
}

% ---- Header ----
\pagestyle{fancy}
\fancyhf{}
\fancyhead[L]{\textit{Lesson 7: Hilbert Spaces and Operators}}
\fancyhead[R]{\thepage}
\renewcommand{\headrulewidth}{0.4pt}

% ---- Theorem Environments ----
\theoremstyle{definition}
\newtheorem{definition}{Definition}[section]
\newtheorem{example}{Example}[section]
\newtheorem{exercise}{Exercise}[section]

\theoremstyle{plain}
\newtheorem{theorem}{Theorem}[section]
\newtheorem{lemma}[theorem]{Lemma}
\newtheorem{proposition}[theorem]{Proposition}
\newtheorem{corollary}[theorem]{Corollary}

\theoremstyle{remark}
\newtheorem{remark}{Remark}[section]

% ---- Colored Boxes ----
\tcbuselibrary{theorems,skins,breakable}

\newtcolorbox{keyresult}[1][]{
    colback=green!5!white,
    colframe=green!50!black,
    fonttitle=\bfseries,
    title=#1,
    breakable
}

\newtcolorbox{rlconnection}[1][]{
    colback=blue!5!white,
    colframe=blue!50!black,
    fonttitle=\bfseries,
    title=#1,
    breakable
}

\newtcolorbox{warning}[1][]{
    colback=red!5!white,
    colframe=red!50!black,
    fonttitle=\bfseries,
    title=#1,
    breakable
}

% ---- Operators ----
\newcommand{\R}{\mathbb{R}}
\newcommand{\N}{\mathbb{N}}
\newcommand{\Q}{\mathbb{Q}}
\newcommand{\Z}{\mathbb{Z}}
\newcommand{\C}{\mathbb{C}}
\newcommand{\Prob}{\mathbb{P}}
\newcommand{\E}{\mathbb{E}}
\newcommand{\indicator}{\mathbf{1}}
\DeclareMathOperator{\interior}{int}
\DeclareMathOperator{\cl}{cl}
\DeclareMathOperator{\spn}{span}
\DeclareMathOperator{\conv}{conv}
\DeclareMathOperator{\dom}{dom}
\DeclareMathOperator{\range}{range}
\DeclareMathOperator{\nullspace}{null}
\DeclareMathOperator{\diam}{diam}

% ---- Title ----
\title{Lesson 21: Fundamental Theorems of Functional Analysis}
\author{Curriculum: A Rigorous Learning Path to Multi-Agent Deep RL}
\date{}

\begin{document}
\maketitle
\tableofcontents
\newpage

\setcounter{section}{6}

\section{Fundamental Theorems of Functional Analysis}\label{sec:fundamental}

We now present the four pillars of linear functional analysis. These theorems
hold in general Banach spaces (not just Hilbert spaces).

\subsection{Hahn--Banach Theorem}

\begin{definition}[Sublinear Functional]\label{def:sublinear}
A function $p \colon X \to \R$ on a real vector space $X$ is \emph{sublinear} if:
\begin{enumerate}[label=(\roman*)]
    \item $p(\alpha x) = \alpha\, p(x)$ for all $\alpha > 0$ and $x \in X$
          (positive homogeneity).
    \item $p(x + y) \leq p(x) + p(y)$ for all $x, y \in X$ (subadditivity).
\end{enumerate}
\end{definition}

\begin{theorem}[Hahn--Banach Theorem, Real Case]\label{thm:hahn-banach}
Let $X$ be a real vector space, $p \colon X \to \R$ a sublinear functional,
and $Y \subset X$ a subspace. If $f \colon Y \to \R$ is a linear functional
with $f(y) \leq p(y)$ for all $y \in Y$, then there exists a linear functional
$F \colon X \to \R$ extending $f$ (i.e., $F|_Y = f$) such that
$F(x) \leq p(x)$ for all $x \in X$.
\end{theorem}

\begin{proof}
\textbf{Step 1: Extension by one dimension.}
Let $x_0 \in X \setminus Y$ and $Y_1 = Y + \R x_0 = \{y + t x_0 : y \in Y, t \in \R\}$.
We seek $c = F(x_0) \in \R$ such that $F(y + t x_0) = f(y) + tc$ satisfies
$F \leq p$ on $Y_1$.

For $t > 0$: $f(y) + tc \leq p(y + tx_0)$, i.e.,
$c \leq \frac{1}{t}[p(y + tx_0) - f(y)] = p(y/t + x_0) - f(y/t)$.
Since $y/t$ ranges over $Y$, we need
$c \leq \inf_{u \in Y}[p(u + x_0) - f(u)]$.

For $t < 0$: setting $s = -t > 0$, we need $f(y) - sc \leq p(y - sx_0)$,
i.e., $c \geq f(y/s) - p(y/s - x_0)$, so
$c \geq \sup_{v \in Y}[f(v) - p(v - x_0)]$.

We must verify:
$f(v) - p(v - x_0) \leq p(u + x_0) - f(u)$ for all $u, v \in Y$.
This reduces to $f(u) + f(v) \leq p(u + x_0) + p(v - x_0)$.
Indeed, $f(u) + f(v) = f(u + v) \leq p(u + v) = p((u + x_0) + (v - x_0))
\leq p(u + x_0) + p(v - x_0)$.

So we can choose $c$ in the interval
$[\sup_v(f(v) - p(v - x_0)),\, \inf_u(p(u + x_0) - f(u))]$.

\textbf{Step 2: Zorn's Lemma.}
Let $\mathcal{F}$ be the set of all pairs $(Z, g)$ where $Y \subset Z \subset X$
is a subspace and $g \colon Z \to \R$ is linear with $g|_Y = f$ and
$g(z) \leq p(z)$ for all $z \in Z$. Order $\mathcal{F}$ by extension:
$(Z_1, g_1) \leq (Z_2, g_2)$ if $Z_1 \subset Z_2$ and $g_2|_{Z_1} = g_1$.

Every chain has an upper bound (take the union of subspaces and the common
extension). By Zorn's Lemma, $\mathcal{F}$ has a maximal element $(Z^*, F)$.
If $Z^* \neq X$, Step~1 would produce a proper extension, contradicting
maximality. Hence $Z^* = X$.
\end{proof}

\begin{corollary}[Extension of Bounded Linear Functionals]\label{cor:hahn-banach-normed}
Let $Y$ be a subspace of a normed space $X$ and $f \in Y^*$. Then there exists
$F \in X^*$ with $F|_Y = f$ and $\|F\|_{X^*} = \|f\|_{Y^*}$.
\end{corollary}

\begin{proof}
Apply Theorem~\ref{thm:hahn-banach} with $p(x) = \|f\| \cdot \|x\|$.
Then $f(y) \leq |f(y)| \leq \|f\| \|y\| = p(y)$, and the extension $F$
satisfies $F(x) \leq \|f\| \|x\|$ and $F(-x) \leq \|f\| \|x\|$,
giving $|F(x)| \leq \|f\| \|x\|$, so $\|F\| \leq \|f\|$.
Since $F$ extends $f$, $\|F\| \geq \|f\|$.
\end{proof}

\begin{corollary}\label{cor:hahn-banach-separation}
For any $x_0 \neq 0$ in a normed space $X$, there exists $f \in X^*$ with
$\|f\| = 1$ and $f(x_0) = \|x_0\|$.
\end{corollary}

\begin{proof}
Define $g \colon \R x_0 \to \R$ by $g(\alpha x_0) = \alpha \|x_0\|$.
Then $\|g\| = 1$ and $g(x_0) = \|x_0\|$.
Extend by Corollary~\ref{cor:hahn-banach-normed}.
\end{proof}

\subsection{Baire Category Theorem and Uniform Boundedness}

We need a preliminary result.

\begin{theorem}[Baire Category Theorem]\label{thm:baire}
Let $(X, d)$ be a complete metric space. If $X = \bigcup_{n=1}^\infty F_n$
where each $F_n$ is closed, then at least one $F_n$ has nonempty interior.
Equivalently, a countable intersection of open dense sets is dense.
\end{theorem}

\begin{proof}
Suppose for contradiction that every $F_n$ has empty interior. Then each
$G_n = X \setminus F_n$ is open and dense. We construct a nested sequence of
closed balls: since $G_1$ is dense, pick $x_1 \in G_1$ and $r_1 < 1$ with
$\overline{B(x_1, r_1)} \subset G_1$. Since $G_2$ is dense and
$B(x_1, r_1)$ is open, pick $x_2 \in G_2 \cap B(x_1, r_1)$ and $r_2 < 1/2$
with $\overline{B(x_2, r_2)} \subset G_2 \cap B(x_1, r_1)$.

Inductively, we obtain $\overline{B(x_{n+1}, r_{n+1})} \subset G_{n+1}
\cap B(x_n, r_n)$ with $r_n < 1/n$. The centers $(x_n)$ form a Cauchy
sequence (since $d(x_m, x_n) < 2/\min(m,n)$ for $m, n$ large), converging
to some $x^* \in X$ by completeness. Then $x^* \in \overline{B(x_n, r_n)}
\subset G_n$ for all $n$, so $x^* \in \bigcap G_n = X \setminus \bigcup F_n
= \emptyset$, a contradiction.
\end{proof}

\begin{theorem}[Uniform Boundedness Principle / Banach--Steinhaus]\label{thm:ubp}
Let $X$ be a Banach space, $Y$ a normed space, and
$\{T_\alpha\}_{\alpha \in A} \subset \mathcal{B}(X, Y)$ a family of bounded
linear operators. If $\sup_\alpha \|T_\alpha x\| < \infty$ for every $x \in X$
(pointwise bounded), then $\sup_\alpha \|T_\alpha\| < \infty$ (uniformly bounded).
\end{theorem}

\begin{proof}
For each $n \in \N$, define the closed set
\[
    F_n = \{x \in X : \|T_\alpha x\| \leq n \text{ for all } \alpha \in A\}
    = \bigcap_{\alpha \in A} \{x : \|T_\alpha x\| \leq n\}.
\]
Each $F_n$ is closed (as an intersection of closed sets, since each
$x \mapsto \|T_\alpha x\|$ is continuous). Since the family is pointwise
bounded, for each $x$ there exists $n$ with $x \in F_n$, so
$X = \bigcup_{n=1}^\infty F_n$.

By the Baire Category Theorem~\ref{thm:baire}, some $F_N$ has nonempty
interior: there exist $x_0 \in X$ and $r > 0$ with $B(x_0, r) \subset F_N$.

For any $x \in X$ with $\|x\| \leq 1$, we have $x_0 + rx \in B(x_0, r) \subset F_N$
and $x_0 \in F_N$, so for all $\alpha$:
\[
    \|T_\alpha(rx)\| = \|T_\alpha(x_0 + rx) - T_\alpha(x_0)\|
    \leq \|T_\alpha(x_0 + rx)\| + \|T_\alpha(x_0)\| \leq N + N = 2N.
\]
Thus $\|T_\alpha x\| \leq 2N/r$ for all $\|x\| \leq 1$, giving
$\|T_\alpha\| \leq 2N/r$ for all $\alpha$.
\end{proof}

\begin{rlconnection}[Uniform Boundedness in Approximation Theory]
The Uniform Boundedness Principle is used to prove negative results in
approximation theory: certain approximation schemes \emph{must} diverge
for some continuous functions. In RL, it guarantees that if a sequence
of operators (e.g., truncated Bellman operators or finite-sample
approximations) converges pointwise, the operator norms remain controlled.
This is crucial for interchanging limits in convergence proofs.
\end{rlconnection}

\subsection{Open Mapping Theorem}

\begin{theorem}[Open Mapping Theorem]\label{thm:open-mapping}
Let $X$ and $Y$ be Banach spaces and $T \colon X \to Y$ a bounded linear
operator that is surjective. Then $T$ is an open map: $T(U)$ is open in $Y$
whenever $U$ is open in $X$.
\end{theorem}

\begin{proof}
\textbf{Step 1:} We show $T(B_X(0, 1))$ contains a ball centered at 0 in $Y$.
Write $B_X = B_X(0,1)$ and $B_Y = B_Y(0,1)$.

Since $T$ is surjective, $Y = \bigcup_{n=1}^\infty T(nB_X) = \bigcup_{n=1}^\infty n \cdot T(B_X)$.
By Baire's theorem, some $\overline{n \cdot T(B_X)}$ has nonempty interior.
By scaling, $\overline{T(B_X)}$ has nonempty interior: there exist
$y_0 \in Y$ and $s > 0$ with $B_Y(y_0, s) \subset \overline{T(B_X)}$.

Since $\overline{T(B_X)}$ is symmetric (if $y \in \overline{T(B_X)}$, then
$-y \in \overline{T(B_X)}$) and convex, we get
$B_Y(0, s) \subset \overline{T(B_X)}$: for any $\|y\| < s$,
$y + y_0 \in \overline{T(B_X)}$ and $-y + y_0 \in \overline{T(B_X)}$,
and $y_0 \in \overline{T(B_X)}$, so
$y = (y + y_0)/2 + (-y_0 + y)/2 + y_0 - y_0$... Let us be more careful.
We have $y + y_0$ and $-y + y_0$ both in $\overline{T(B_X)}$, and
$\overline{T(B_X)} - \overline{T(B_X)} \subset \overline{T(2B_X)}$
(since $T(B_X) + T(B_X) \subset T(2B_X)$). In particular,
$(y + y_0) - (-y + y_0) = 2y \in \overline{T(2B_X)}$, so
$y \in \overline{T(B_X)}$. Hence $B_Y(0, s) \subset \overline{T(B_X)}$.

\textbf{Step 2:} We show $\overline{T(B_X)} \subset T(2B_X)$, i.e.,
the closure can be ``opened.'' Take $y \in \overline{T(B_X)}$. From Step 1,
$B_Y(0, s) \subset \overline{T(B_X)}$, so for each $r > 0$,
$B_Y(0, rs) \subset \overline{T(rB_X)}$.

Choose $x_1 \in B_X$ with $\|y - Tx_1\| < s/2$. Then
$y - Tx_1 \in B_Y(0, s/2) \subset \overline{T(\frac{1}{2}B_X)}$, so choose
$x_2 \in \frac{1}{2}B_X$ with $\|y - Tx_1 - Tx_2\| < s/4$.
Inductively, choose $x_n \in 2^{-(n-1)} B_X$ with
$\|y - \sum_{k=1}^n Tx_k\| < s \cdot 2^{-n}$.

Then $\sum_n \|x_n\| < 1 + 1/2 + 1/4 + \cdots = 2$. Since $X$ is complete,
$x = \sum_{n=1}^\infty x_n$ converges with $\|x\| < 2$, and
$Tx = \lim_{n \to \infty} T(\sum_{k=1}^n x_k) = y$. So $y \in T(2B_X)$.

\textbf{Step 3:} Combining Steps 1 and 2:
$B_Y(0, s) \subset \overline{T(B_X)} \subset T(2B_X)$, so
$B_Y(0, s/2) \subset T(B_X)$.

\textbf{Step 4:} $T$ is open. Let $U \subset X$ be open and $y_0 = Tx_0 \in T(U)$.
Choose $r > 0$ with $B_X(x_0, r) \subset U$. Then
$T(U) \supset T(B_X(x_0, r)) = Tx_0 + T(rB_X) \supset y_0 + B_Y(0, rs/2)$,
so $T(U)$ is open.
\end{proof}

\begin{corollary}[Bounded Inverse Theorem]\label{cor:bounded-inverse}
If $T \colon X \to Y$ is a bounded linear bijection between Banach spaces,
then $T^{-1}$ is bounded.
\end{corollary}

\begin{proof}
$T$ is an open map, so $T^{-1}$ is continuous (preimages of open sets under
$T^{-1}$ are images of open sets under $T$).
\end{proof}

\subsection{Closed Graph Theorem}

\begin{definition}[Closed Graph]\label{def:closed-graph}
A linear operator $T \colon X \to Y$ between normed spaces has a
\emph{closed graph} if its graph
$G(T) = \{(x, Tx) : x \in X\} \subset X \times Y$ is closed in the
product topology (equivalently, in the norm $\|(x,y)\| = \|x\| + \|y\|$).
Equivalently: if $x_n \to x$ in $X$ and $Tx_n \to y$ in $Y$, then $y = Tx$.
\end{definition}

\begin{theorem}[Closed Graph Theorem]\label{thm:closed-graph}
Let $X$ and $Y$ be Banach spaces. A linear operator $T \colon X \to Y$ is
bounded if and only if it has a closed graph.
\end{theorem}

\begin{proof}
$(\Rightarrow)$ If $T$ is bounded and $x_n \to x$, $Tx_n \to y$, then
$Tx_n \to Tx$ by continuity, so $y = Tx$ by uniqueness of limits.

$(\Leftarrow)$ Assume $G(T)$ is closed. Consider $X \times Y$ with the norm
$\|(x,y)\| = \|x\| + \|y\|$, which is a Banach space. $G(T)$ is a closed
subspace, hence itself a Banach space.

Define $\pi_1 \colon G(T) \to X$ by $\pi_1(x, Tx) = x$ and
$\pi_2 \colon G(T) \to Y$ by $\pi_2(x, Tx) = Tx$.
Then $\pi_1$ is bounded, linear, and bijective. By the Bounded Inverse
Theorem (Corollary~\ref{cor:bounded-inverse}), $\pi_1^{-1} \colon X \to G(T)$
is bounded. Since $T = \pi_2 \circ \pi_1^{-1}$ and both $\pi_2$ and
$\pi_1^{-1}$ are bounded, $T$ is bounded.
\end{proof}

\begin{remark}\label{rem:closed-graph-usage}
The Closed Graph Theorem is often the most convenient way to prove that
an operator is bounded: instead of estimating $\|Tx\|$ directly, one
only needs to verify the sequential closedness condition, which is often
easier in applications.
\end{remark}


%% ============================================================================
%% SECTION 8: FIXED POINT THEOREMS
%% ============================================================================


\end{document}